% Generated draft: requires mathematical and source review.
\documentclass[11pt]{article}
\usepackage[margin=25mm]{geometry}
\usepackage{fontspec}
\setmainfont{Latin Modern Roman}
\usepackage{amsmath,amssymb,mathtools,unicode-math}
\setmathfont{Latin Modern Math}
\usepackage{enumitem,needspace,xurl,hyperref}
\setlength{\parindent}{0pt}
\setlength{\parskip}{5pt}
\setlength{\emergencystretch}{3em}
\newcommand{\N}{\mathbb{N}}
\newcommand{\Z}{\mathbb{Z}}
\newcommand{\Q}{\mathbb{Q}}
\newcommand{\R}{\mathbb{R}}
\newcommand{\C}{\mathbb{C}}
\newcommand{\F}{\mathbb{F}}
\newcommand{\A}{\mathbb{A}}
\newcommand{\PP}{\mathbb{P}}
\newcommand{\E}{\mathbb{E}}
\newcommand{\eps}{\varepsilon}
\newcommand{\dd}{\,\mathrm{d}}
\begin{document}
\section*{UnsolvedMath ID: 30001801}
\textbf{Source problem code:} catalogue entry
\subsection{Definitions and mathematical statement}
A group is a virtual Poincar\'e-duality group if it has a torsion-free subgroup of finite index that is a Poincar\'e-duality group. The latter means a group satisfying group-cohomological Poincar\'e duality, with its orientation module allowed to be nontrivial. A finite Poincar\'e complex of dimension \(d\) is a connected finite CW complex equipped with an orientation local system and a fundamental class inducing Poincar\'e duality with local coefficients. Write \(\widetilde X\) for its universal cover and \(\simeq\) for homotopy equivalence.

The Spivak normal fibration \(\nu_X\) is the stable spherical fibration associated with the Poincar\'e duality of \(X\). Reducibility requires specifying a bundle category: a vector-bundle reduction is a lift of its classifying map \(X\to BG\) through \(BO\to BG\), where \(BG\) classifies stable spherical fibrations. Piecewise-linear and topological bundle reductions instead use \(BPL\to BG\) and \(BTOP\to BG\), respectively. The supplied record does not specify which reduction is intended.

\textbf{Supplied formulation.} Let \(\Gamma\) be a virtual Poincar\'e-duality group described in the record as having cohomological dimension \(k\), and suppose its Farrell cohomology is periodic of period \(n+1\). Let \(X\) be a finite Poincar\'e complex supplied by the referenced construction, satisfying
\[
\dim X=n+k,\qquad \pi_1(X)\cong\Gamma,\qquad \widetilde X\simeq S^n.
\]
The two questions in the clean statement are whether \(\nu_X\) is reducible in the intended bundle category, and whether there exists a closed manifold \(M\) such that
\[
X\simeq M.
\]
Here a closed manifold is compact and without boundary; its intended category is not specified.

\textbf{Unresolved quantification and hypotheses.} The original wording asks whether \(X\) \emph{can be constructed} with reducible Spivak normal fibre space, and whether \emph{such an} \(X\) can be homotopy equivalent to a closed manifold. This is an existence question about a choice of \(X\) supplied by a theorem. The clean wording instead asks about an already chosen resulting \(X\). These do not specify the same quantification over possible complexes, so the draft retains both readings. The construction theorem itself is absent, and its full hypotheses and permitted choices cannot be recovered from this record.

The phrase ``cohomological dimension \(k\)'' also requires clarification: ordinary integral cohomological dimension and virtual cohomological dimension need not agree for a virtual Poincar\'e-duality group. The record does not specify which is meant. Likewise, it does not state the coefficient convention for periodic Farrell cohomology, whether \(n+1\) is required to be the least positive period, or the parameter restrictions imposed by the referenced theorem. No additional hypotheses or choices of these conventions are asserted here.
\subsection{Short English statement}
For a finite Poincar\'e complex \(X\) arising from the stated group-theoretic construction, ask whether its Spivak normal fibration admits the intended bundle reduction and whether \(X\) is homotopy equivalent to a closed manifold. The original source wording asks whether a suitable \(X\) can be chosen; the clean statement asks about a resulting \(X\). The missing construction theorem, dimension convention, and bundle and manifold categories require clarification.
\subsection{Sources}
\begin{itemize}
\item Supplied local UnsolvedMath registry record, permanent ID \(30001801\), exact source problem code \texttt{catalogue entry}, title ``Reducibility and Manifold Realization of Finite Poincar\'e Complexes.'' The record belongs to the Oberwolfach Reports Open Problems collection and gives proposed year 2011; no individual proposer is supplied.
\item Supplied report citation: \emph{New Directions in Algebraic \(K\)-Theory} (2011), Oberwolfach Reports, pp.\ 1469--1509, DOI \texttt{10.4171/owr/2011/27}. Source URL and literature-source link labelled ``original questions'': \url{https://doi.org/10.4171/owr/2011/27}. The record describes Oberwolfach Reports as documenting mathematical workshops and their open-problem discussions; no precise page or theorem number for this entry is supplied.
\item Supplied literature-source link labelled ``Hambleton'': \url{https://ms.mcmaster.ca/ian/papers/hambleton_2019.pdf}. The record attributes to Hambleton a reducibility result for the Spivak normal fibre space under substantial finite index hypotheses. It supplies neither those hypotheses nor a full bibliographic citation or theorem number, so no precise theorem is asserted here.
\item Literature assessment checked at \texttt{2026-08-21}: the supplied curation labels the problem ``partially solved'' and states that the two general surgery questions in the extracted entry are not fully settled. This is supplied dated curation, not an independently verified literature assessment.
\item The supplied registry labels the statement \texttt{corrected\_verified} and reports that it was checked against the official Oberwolfach report. This draft preserves that verification claim as source metadata; it does not independently verify the report or resolve the formulation issues identified above.
\end{itemize}
\end{document}
